
%%%%%%%%%%%%%%%%%%%%%%%%%%%%%%%%%%%%%%%%%%%%%%%%%%%%%%%%%%%%%%%%%%%%
%%%%%%%% Released for translation     10.4.2023 %%%%%%%%%%%%%%%%%%%%
%%%%%%%% Re-released for translation  22.6.2023 %%%%%%%%%%%%%%%%%%%%
%%%%%%%% Re-released for translation 26.12.2023 %%%%%%%%%%%%%%%%%%%%
%%%%%%%% Re-released for translation 04.01.2024 %%%%%%%%%%%%%%%%%%%%
%%%%%%%%%%%%%%%%%%%%%%%%%%%%%%%%%%%%%%%%%%%%%%%%%%%%%%%%%%%%%%%%%%%%

\noindent
\brcolor{\LARGE Preface}\label{primer-preface}\label{preface} 
\addcontentsline{toc}{section}{\brcolor{Preface}}
\minitoc
 
\begin{flushright}\label{front-Dogma}
\sort{The Triptych Dogma}\\[4mm]
\addcontentsline{toc}{subsection}{\bbcolor{The Triptych Dogma}}
\sf In order to \textsl{specify} \bmcolor{software},\\
\sf we must understand its requirements.\\[1mm]
\sf In order to \textsl{prescribe} \bmcolor{requirements}\ysfchg{,}\\
\sf we must understand the \bmcolor{domain}.\\[1mm]
\sf So we must \bbcolor{study, analyse} and \bbcolor{describe} domains.
\end{flushright}\rm%
\index{pdefind}{The Triptych Dogma}\index{pdefind}{Triptych!Dogma, The}%\index{pdefind}{Dogma, The Triptych}

\subsection*{\bbcolor{Domains -- What Are They\,?}}
\addcontentsline{toc}{subsection}{\bbcolor{Domains -- What Are They\,?}}

\pos{\normalsize}{\HHHH}

\label{1stdd}
\ysf{\domaindefinition}
\pos{\psno}{\mnewfoil}

\renewcommand{\newdbsquare}{\dbsquare}

\begynd
\pind \sort{Examples} \ysfchg{o}f domains are:
\begynd
\pind {\sl rail, road, sea and air transport;
\pind water, oil and gas pipelines;
\pind industrial manufacturing;
\pind the financial service industry: clients, banks, credit cards, stocks, etc.;
\pind consumer, retail and wholesale markets;
\pind health care;}
\pind et cetera.
\afslut
\afslut
\subsection*{\bbcolor{Aim and Objectives}}\addcontentsline{toc}{subsection}{\bbcolor{Aim and Objectives}}  

\vspace*{-4mm}

\begin{itemize}
\item The \sort{aim} of this \monograph\ is to contribute to a 
      methodology for analysing and describing domains.
  
\item The \sort{objectives} -- in the sense of `how is the aim
      achieved' -- is reflected in the structure and contents
      and the didactic approach of this \monograph.
\item The main elements of \ysfchg{our} approach -- along one
  concept-axis -- can be itemized: 
\begin{itemize}
\item There is the founding of our analysis \& description
      approach in providing a base \sort{philosophy},
      cf.\,Chapter\,\ref{chap2.tex.Philosophy}. 
\item There is the application of ideas of \sort{taxonomy} 
      to understand the possibly hierarchical structuring
      of domain phenomena respectively the understanding of properties
      of phenomena and relations between them.
\item There {are} the notions \sort{endurants} and \sort{perdurants} --
      with \sfsl{endurants}  \index{pdefind}{endurant} being the phenomena
\begynd
\pind that can be observed, or conceived and described, as \nyl a ``complete
      thing'' at no matter which given snapshot of time\pos{
      \cite[Vol.\,I, pg.\,656]{OED}}{},
\afslut
      and \sfsl{perdurants}  \index{pdefind}{perdurant} being the phenomena 
\begynd
\pind for which only a fragment exists\pos{}{\\} if we look at or
      touch them\pos{}{\\} at
      any given snapshot in time\pos{
      \cite[Vol.\,II, pg.\,1552]{OED}}{}.
\afslut
\item There is the introduction of base elements of
      \sort{calculi}  \index{pdefind}{calculi} for analysing and
      describing domains. \index{pdefind}{calculus!analysis}
\item There is the application of ideas of \sort{ontology} 
      to understand the possibly hierarchical structuring
      of these calculi.\index{pdefind}{ontology} 
\item And finally there is the notion of
      \index{pdefind}{transcendental deduction} \sort{transcendental
      deduction}, cf.\,Sect.\,\ref{sec:Transcendence}, for ``morphing''
      certain kinds of endurants into certain kinds of perdurants,
      Chapter \ref{chap6.tex.1}. 
\end{itemize}
\item Along another conceptual-axis the below are further  elements
        of our approach:
\begin{itemize}
\item We consider domain descriptions, requirements prescriptions and
      software design specifications to be \brcolor{mathematical}
      \index{pdefind}{mathematics} quantities.
\item And we consider them basically in the sense of
      \sort{recursive function theory}
      \index{pdefind}{recursive!function theory} \cite[Hartley Rogers,
      1952]{Rogers67} and 
      \sort{type theory}  \index{pdefind}{type!theory} \cite[Benjamin
      Pierce, 1997]{pierce1997}.  
\end{itemize}
\end{itemize}

\vspace*{-4mm}

\subsection*{\bbcolor{Methodology}}\addcontentsline{toc}{subsection}{\bbcolor{Methodology}}\label{front:Methodology}

\begynd
\pind By a \sort{method} we shall understand
\begynd
\pind a set of \sort{principles}\index{pdefind}{principle}%
\footnote{By a \sort{principle} we mean:
  \sfsl{a principle is a proposition or value that is a guide for
    behavior or evaluation \wiki, i.e., code of conduct}}
and \sort{procedures}\index{pdefind}{procedure}%
\footnote{By a
  \sort{procedure} we mean: \sfsl{instructions or recipes, a set of commands
    that show how to achieve some result, such as to prepare or make
    something \wiki, i.e., an established way of doing something}} 
\pind for selecting and applying a set of
\sort{techniques}\index{pdefind}{technique}%
\footnote{By a \sort{technique} we mean: \sfsl{a technique, or skill, is
    the learned ability to perform an action with determined results
    with good execution often within a given amount of time, energy,
    or both \wiki, i.e., a way of
    carrying out a particular task}}
and \sort{tools}\index{pdefind}{tool}%
\footnote{By a
  \sort{tool} we mean: \sfsl{a tool is an object that can extend an
    individual's ability to modify features of the surrounding
    environment \wiki}} to a problem  
\afslut
\pind in order to achieve an orderly construction of a
      \sort{solution}, i.e., an \sort{artefact}.
\afslut
\mnewfoil

\begynd
\pind By \sort{methodology} we shall understand 
\begynd
\pind the \sfsl{study \& application} of one or more methods.
\afslut
\afslut
\mnewfoil

\pind By a \sort{formal method} we shall understand a method
\begin{itemize}
\item  whose \sfsl{principles} include that of considering its artefacts as
      \sfsl{mathematical} quantities, of \sfsl{abstraction}, etc.;
\item  whose decisive \sfsl{procedures} include that of
\begin{itemize}
\item the sequential
      analysis \ysfchg{\&} description of first endurants, then perdurants,
      and,
\item  within the analysis  \ysfchg{\&} description of endurants, the
      sequential  analysis   \ysfchg{\&} description of first their external
      qualities and then their internal qualities,
\item etc.; 
\end{itemize}
\item  whose \sfsl{techniques} include those of specific ways of
       specifying properties; and
\item  whose \sfsl{tools} include those of one or more
      \sort{formal languages.} 
\end{itemize}
\mnewfoil

\noindent
\begynd
\pind By a \sort{language} we shall here understand \nyl a set of
strings of characters, i.e., sentences, 
\begynd
\pind sentences which are structured according to some \sort{syntax},
      i.e., \sort{grammar}, 
\pind are given meaning by some \sort{semantics},
\pind and are used according to some \sort{pragmatics}. 
\afslut
\afslut

\begynd
\pind By a \sort{formal language} we shall here understand a language 
\begynd
\pind whose \sfsl{syntax} and \sfsl{semantics} can both be expressed
      \sort{mathematically} 
\pind and for whose sentences one can \sort{rationally reason}
      (\sfsl{argue, prove}) \sort{properties}. 
\afslut

\pind We refer to Chapter\,1 of \cite{BjornerMonograph2020} for an 8
      page, approximately 50 entries set of concept definitions such
      as the above.
      
\pind We refer to the \textsf{Method} index, Sect.\,\vref{label.dadmethod}.
\afslut
      
\treprikker

\noindent
\pind In this \primer\ we shall use
\begynd 
\pind the formal specification language, \texttt{RSL}, 
\pind the \texttt{R}AISE\footnote{\sort{RAISE: R}igorous
      \sort{A}pproach\ysfchgii{ }
      to \ysfchg{\sort{I}ndustrial } \sort{S}oftware \sort{E}ngineering,
      \cite{RaiseMethod}} \texttt{S}pecification 
      \texttt{L}anguage, \cite{RSL} -- 
\pind and we shall
      notably rely on \texttt{RSL}'s adaptation of \sort{CSP}, Tony
      Hoare's \sfsl{Communicating Sequential Processes} \citecsp;
\pind and we shall propagate a definitive method for the study and
      description of domains.
\afslut
\afslut

\subsection*{\bbcolor{An Emphasis}}\addcontentsline{toc}{subsection}{\bbcolor{An Emphasis}} 

\anemphasis{Domains exhibit endurants and perdurants. A domain model
  is therefore something that defines the \sfsl{nouns} (roughly speaking the
  endurants) and \sfsl{verbs} (roughly speaking the {perdurants}) -- and
  their combination -- of a
  \sfsl{language} spoken and used in writing by the practitioners of the domain. Not an
  instantiation of nouns, verbs and their combination,
  but all possible and sensible instantiations.}{}

\subsection*{\bbcolor{A Caveat}}\addcontentsline{toc}{subsection}{\bbcolor{A Caveat}}

Experienced \texttt{RSL} \cite{RSL} readers might observe our,
perhaps cavalier (offhand), use of  \texttt{RSL}. Perhaps, in some places, the 
syntax of  \texttt{RSL} clauses is not quite right. Our non-use of
\texttt{RSL}'s module (\texttt{Scheme}\ysfchg{, } \texttt{Class} and \texttt{Object})
constructs force \ysfchg{us } to declare \textsf{channel}s in the same way
\textsf{type}s, \textsf{value}s and \textsf{variable}s are introduced.

\begin{itemize}
%%\item[] \hfill \epsfig{file=db.eps,height=25mm}

%%\vspace*{-12mm}
\item[] \hfill   Dines Bj{\o}rner \& Yang ShaoFa

\item[] \hfill \todaytime
\end{itemize}


\label{preface.n}
%%  LocalWords:  analyse defind artefactual et cetera analysing RSL
%%  LocalWords:  endurants perdurants endurant perdurant Hartley AISE
%%  LocalWords:  pdefind artefact artefacts dadmethod igorous pproach
%%  LocalWords:  oftware ngineering pecification anguage CSP Hoare's
%%  LocalWords:  instantiation instantiations ndustrial Bj rner
%%  LocalWords:  ShaoFa
